\documentclass[../main.tex]{subfiles}
\begin{document}
%==========================================TIÊU ĐỀ TẬP========================================
\begin{table}[h]
    \begin{adjustwidth}{-1cm}{}
        \begin{tabular}{>{\centering\arraybackslash}p{6cm}|>{\raggedright\arraybackslash}p{12.5cm}}
            \multirow{5}{6cm}
            % Cột bên trái: ÉPISODE và số tập
            {\\[-40pt]
                \flaregothic\fontsize{20pt}{20pt}\selectfont \centering
                \color{eptype}HISTOIRE\color{black}\\
                \freshbot\fontsize{72pt}{72pt}\selectfont
                \color{epnum}10\\[6pt]
                \cafeteria\fontsize{20pt}{20pt}\selectfont\color{epnum}(D-10)\color{black}
                \\[18pt]
                \color{black}
            }
             &
            % Dòng thứ nhất: tên tập tiếng Nhật
            {\fotskip\fontsize{18pt}{18pt}\selectfont \ruby{小}{ちい}さな\ruby{手}{て}が\ruby{紡}{つむ}ぐ\ruby{温}{ぬく}もり}
            \\[6pt]
            % Dòng thứ hai: tên tập tiếng Anh
             & {\cafeteria\fontsize{18pt}{18pt}\selectfont The Gentle Anchor of Naomi}                                                                                    \\[4pt]
            % Dòng thứ ba: tên tập tiếng Việt
             & {\vnmsans\fontsize{18pt}{18pt}\selectfont Trái tim của ký túc xá}                                                                                          \\[8pt]
            % Dòng thứ tư: Nhân vật xuất hiện
             & {\caviar{\fontsize{14pt}{14pt}\selectfont Nhân vật: \emp{kuon2} \emp{ryuuhou2} \emp{kaigou2} \emp{suzuno2} \emp{naomi2} \emp{game} \emp{coco} \emp{aloe}}}
            \\[8pt]
            % Dòng cuối cùng: Release Date
            % & {\continuum\fontsize{16pt}{16pt}\selectfont Release Date: 07/02/2021}\\[18pt]
        \end{tabular}
    \end{adjustwidth}
\end{table}
%=========================================TRUYỆN CHÍNH========================================
\color{black}

\txtbx[2]
{{\color{vol} \uline{Tóm tắt:}} Nhận thấy Coco và Aloe vẫn luôn mang những nét trầm tư trĩu nặng từ quá khứ, cô bé út Naomi quyết định bí mật chuẩn bị một bữa tiệc trà nhỏ cùng những món quà thủ công để cảm ơn hai chị. Bằng sự hồn nhiên, khả năng thấu thị xoa dịu tâm hồn, nét vẽ ngây thơ cùng một buổi chiều dạo chơi ấm áp, Naomi đã làm mềm lòng cô rồng gai góc Coco và tiếp thêm ngọn lửa âm nhạc cho Aloe. Bữa tiệc trà ngập tràn tiếng cười khẳng định Naomi chính là mỏ neo cảm xúc, biến ký túc xá Hikawa từ chốn trọ tạm bợ thành một mái ấm gia đình thực thụ.}

Ánh nắng dịu dàng của một buổi chiều đầu thu buông xuống hiên nhà Hikawa, rọi qua kẽ lá cây phong trước sân tạo nên những vệt sáng lấp lánh trên sàn gỗ. Căn nhà gỗ hai tầng vốn luôn náo nhiệt những ngày qua bỗng rơi vào một khoảng lặng êm ả hiếm hoi. Kuon và Kaigou vẫn đang bận rộn với công việc quản lý và kỹ thuật tại văn phòng, trong khi Ryuuhou tranh thủ ngủ bù sau một buổi tối chạy show underground dồn dập.

Ở góc hành lang cầu thang nhìn ra ngoài cảnh vật nườm nượp đường phố, Naomi đang ngồi ngay ngắn trên bậc thềm gỗ, đôi chân nhỏ nhắn khẽ đung đưa trong không trung. Đôi mắt màu xanh lục bảo trong vắt của cô bé - thứ ánh nhìn tĩnh lặng như mặt hồ rừng sâu đã trải qua hơn một thiên niên kỷ - đang lặng lẽ quan sát từng chuyển động nhỏ trong ngôi nhà.

Dù mang hình hài của một đứa trẻ tám tuổi sau khi được tái sinh nhờ sự cộng hưởng linh hồn với Suzuno và sự can thiệp mã nguồn từ Game, Naomi sở hữu một khả năng thấu thị kỳ diệu. Cô bé có thể nhìn thấy những dòng chảy năng lượng mờ ảo bao bọc quanh mỗi con người.

Và lúc này, ở hai góc khác nhau của ký túc xá, cô em út nhà Hikawa nhìn thấy hai luồng ánh sáng mang những vết rạn nứt sâu thẳm.

Ở góc ban công phòng khách chính tầng năm, Coco đang tựa lưng vào lan can, ngón tay xoay xoay một chiếc bật lửa kim loại chưa châm. Nàng Rồng huyền thoại với mái tóc cam rực rỡ và dáng vẻ ngạo nghễ thường ngày bỗng toát lên một nỗi cô đơn trầm mặc. Dù luôn tỏ ra mạnh mẽ và sẵn sàng bảo vệ mọi người, quanh người Coco vẫn vương vấn những dải năng lượng màu xám bạc - tàn dư của những cơn sóng gió dư luận khủng khiếp mà cô từng phải một mình gánh chịu trước khi quyết định rời đi.

Còn phía bên kia hành lang, tại tầng hai, sau cánh cửa phòng hé mở, Aloe đang ngồi bó gối bên chiếc bàn học. Trước mặt cô gái quỷ nhỏ là những khuông nhạc viết dở bị gạch xóa chằng chịt và chiếc bút chì đã bị cắn mòn đuôi. Luồng khí tức quanh Aloe lại là một màu tím nhạt e ấp, co cụm lại như một chú nhím nhỏ sợ hãi thế giới bên ngoài sau cú ngã đau đớn thời điểm mới bước chân vào nghề.

Cả hai người họ đều được Kuon đưa về che chở, nhưng dường như trong thâm tâm, họ vẫn chỉ dám coi mình là những ``kẻ ở trọ tạm thời'', luôn giữ một khoảng cách vô hình vì mặc cảm quá khứ.

Naomi khẽ chớp mắt, đôi con ngươi lục bảo ánh lên một tia sáng ấm áp. Cô bé quay đầu lại nhìn vào trong bếp, nơi Suzuno đang tỉ mẩn xếp từng chiếc tách gốm ra khay.

``Suzuno-oneechan,'' cô bé cất giọng trong trẻo, bước chân nhẹ nhàng không phát ra tiếng động tiến lại gần cô chị gái thứ.

Suzuno ngẩng lên, đôi mắt hai màu đỏ - xanh lam lấp lánh vẻ dịu dàng. Cô mỉm cười xoa nhẹ mái tóc mềm của em gái út: ``Sao thế Naomi-chan? Em đói bụng rồi à? Chị đang chuẩn bị chút bánh quy nướng đây.''

``Không phải đâu ạ,'' Naomi lắc đầu, hai bàn tay nhỏ đan vào nhau. ``Naomi muốn làm một điều gì đó cho Coco-oneechan và Aloe-oneechan. Hai chị ấy\dots\ trong lòng vẫn còn nhiều vệt buồn lắm. Naomi muốn làm một bữa tiệc trà nhỏ, và tự tay làm quà tặng các chị ấy để cảm ơn. Chị giúp em được không?''

Nghe những lời chân thành phát ra từ tâm hồn thuần khiết của cô bé, cô chị Ba xúc động mỉm cười rạng rỡ. Tiếng chuông vàng trên chiếc vòng tay đỏ của Suzuno khẽ leng keng theo từng nhịp gật đầu: ``Ý hay lắm đó, Naomi! Chị sẽ phụ trách nướng mẻ bánh quy thơm nhất và pha hồng trà hoa cúc. Còn phần quà bất ngờ, em đã nghĩ ra sẽ làm gì chưa?''

Naomi nghiêng đầu suy nghĩ một lát rồi mỉm cười: ``Em biết mình phải làm gì rồi ạ!''

\ct

Phân đoạn đầu tiên trong kế hoạch nhỏ của Naomi bắt đầu tại căn phòng của Aloe.

Tiếng gõ cửa cộc cộc vang lên khe khẽ. Nàng Succubus giật mình ngẩng đầu, vội vàng lấy tập tài liệu che đi những bản nhạc dang dở trên bàn rồi mới nhỏ giọng: ``A\dots\ cửa không khóa đâu, vào đi.''

Cánh cửa trượt mở ra, để lộ chiếc đầu nhỏ nhắn cùng đôi mắt to tròn của Naomi. Cô bé ôm theo một xấp giấy vẽ dày cộp, hộp bút màu sáp và một chiếc kéo thủ công đầu tròn.

``Aloe-oneechan có bận không ạ?'' cô bé lễ phép hỏi.

Aloe tròn mắt ngạc nhiên, vội vàng đứng dậy kéo một chiếc đệm êm đặt xuống sàn: ``Không, chị không bận gì đâu. Naomi cần chị giúp gì sao?''

Naomi lon ton chạy lại, ngồi bệt xuống đệm rồi ngẩng khuôn mặt bầu bĩnh lên: ``Em muốn vẽ vài tấm thiệp đẹp để tặng mọi người, nhưng Naomi vẽ hoa lá còn vụng lắm. Mọi người bảo Aloe-oneechan có đôi tay rất khéo và hát cũng rất hay. Chị có thể dạy em vẽ và dạy em một bài hát ngắn được không ạ?''

Lời đề nghị bất ngờ khiến cô chị khựng lại. Hai má cô bé quỷ nhỏ ửng hồng, nhưng ngay sau đó ánh mắt lại thoáng vẻ ngập ngừng, tự ti: ``Chị\dots\ chị vẽ cũng bình thường thôi mà. Còn hát\dots\ lâu rồi chị không có hát riêng trước mặt ai cả, sợ dở làm Naomi chê mất\dots\ Kuon-kun với Ryuuhou-chan lúc chị chuyển về cũng bị chê bai đủ kiểu rồi\dots''

``Sẽ không dở đâu ạ!'' Naomi quả quyết, đưa bàn tay nhỏ bé mềm mại nắm lấy những ngón tay đang run nhẹ của Aloe. Một luồng khí tức mát lành, thanh khiết như dòng suối đầu nguồn từ bàn tay Naomi truyền sang, xoa dịu đi sự căng thẳng bấy lâu trong lồng ngực Aloe. ``Giọng của Aloe-oneechan ấm lắm, giống như tiếng gió thổi qua rặng tre mùa xuân vậy. Em rất muốn được nghe.''

Nhìn vào đôi mắt trong veo không một chút gợn đục của cô bé tám tuổi, mọi phòng tuyến phòng thủ và nỗi mặc cảm trong lòng cô nàng Quỷ hoàn toàn sụp đổ. Một nụ cười tự nhiên, rạng rỡ lần đầu tiên nở trên môi cô bé quỷ nhỏ sau chuỗi ngày dài u uất.

``Được rồi\dots\ vậy chị em mình cùng vẽ nhé,'' Aloe ngồi xuống bên cạnh Naomi, cầm lấy cây bút màu sáp màu hồng. ``Đầu tiên, vẽ cánh hoa thì em phải đưa nét cong nhẹ thế này này\dots''

Trong căn phòng nhỏ, từng nét vẽ ngây ngô dần thành hình. Aloe kiên nhẫn chỉ cho Naomi cách phối màu, từ những bông hoa cẩm chướng rực rỡ đến hình vẽ chibi ngộ nghĩnh của cả đại gia đình Hikawa. Khi vẽ xong những tấm thiệp, Aloe khẽ với tay lấy cây đàn ukulele nhỏ treo trên tường.

Ngón tay cô lướt nhẹ trên từng phím đàn, phát ra những âm thanh mộc mạc, trong trẻo. Aloe cất tiếng hát khẽ khàng - một giai điệu đồng dao nhẹ nhàng về những đám mây mùa hạ. Naomi chăm chú lắng nghe rồi cất giọng hòa theo. Hai giọng hát, một non nớt trong veo, một ngọt ngào sâu lắng, quyện vào nhau xua tan đi mọi bóng tối u ám từng bám chặt trong căn phòng này.

\ct

Sau khi hoàn thành bài học vẽ và âm nhạc với Aloe, Naomi thực hiện bước thứ hai của kế hoạch: tiếp cận ``mục tiêu khó nhằn nhất'' - là Kiryu Coco.

Lúc này, Coco đang ngồi khoanh chân trên chiếc ghế dài ở phòng khách, tay cầm điện thoại lướt xem tin tức, gương mặt vẫn mang vẻ lạnh lùng khó gần.

Bỗng nhiên, chiếc vạt áo thun của cô bị giật nhẹ hai cái.

Nàng Rồng cúi đầu xuống, bắt gặp Naomi đang đứng sát bên cạnh, trên đầu đội một chiếc mũ cói rộng vành, hai tay xách theo một chiếc giỏ mây nhỏ.

``Gì thế Naomi?'' Coco nhướng mày, giọng nói có phần trầm khàn bộc trực cố hữu. ``Kuon-paisen với bà chị quái vật của nhóc chưa về đâu. Cần gì thì bảo Suzuno ấy.''

Cô em út nhà Hikawa không hề sợ hãi trước vẻ mặt bặm trợn ấy, ngược lại còn tiến thêm một bước, vươn tay nắm chặt lấy ngón tay út to lớn của cô rồng: ``Naomi không tìm Onii-chan hay Kaigou-oneechan. Naomi muốn rủ Coco-oneechan đi dạo công viên bờ sông với em cơ.''

Nàng Rồng giật mình, mắt mở to nhìn ngón tay út của mình bị bàn tay bé xíu nắm chặt. Cô định rút tay lại theo phản xạ, nhưng cảm giác mềm mại và ấm áp từ bàn tay đứa trẻ khiến cô không nỡ dùng lực.

``Này\dots\ chị mày bận lắm đấy nhá,'' Coco càu nhàu, nhưng ngón tay lại vô thức co lại để giữ lấy tay bé. ``Với lại trời ngoài kia nắng chết cha ra, đi ra đấy làm gì cho mệt xác?''

``Đi với em một chút thôi mà, Coco-oneechan\~{}'' Naomi ngước mắt lên, làm vẻ mặt nũng nịu ngây thơ; một chiêu thức mà dù là rồng thần ngàn năm hay bất kỳ ai cũng không thể chống đỡ nổi.

Thấy rằng bản thân không thể nào cưỡng được sự đáng yêu của cô em út, Coco thở hắt ra một hơi dài bất lực, đưa tay gãi gãi đầu: ``Thôi được rồi! Coi như hôm nay chị rảnh rỗi sinh nông nổi đi dạo với nhóc một vòng. Nhưng cấm có chạy lung tung đấy nhé!''

``Vâng ạ\~{}'' Naomi hồn nhiên đáp lại, lon ton chạy ra khỏi phòng.

Thế là, trên con đường dốc dẫn ra công viên ven sông, người dân xung quanh khu phố của phường Anwa được chứng kiến một cảnh tượng vô cùng thú vị: một cô gái cao ráo, dáng đi ngông nghênh đậm chất giang hồ đang cẩn thận dắt tay một bé gái nhỏ nhắn đáng yêu.

Dọc đường, Coco liên tục lẩm bẩm càm ràm, nhưng mắt thì không rời Naomi nửa bước. Thấy một chiếc xe đạp chạy ẩu phóng qua, nàng Rồng nhanh như cắt bước tới chắn ngang phía ngoài, trừng mắt lườm gã lái xe khiến người kia phải rụt cổ phóng vội.

Đi ngang qua quầy kem bông gòn, thấy Naomi nhìn chằm chằm vào cây kẹo màu hồng khổng lồ, Coco không nói không rằng rút ví mua ngay một cây kẹo to nhất đưa cho cô bé: ``Cầm lấy. Ăn cẩn thận kẻo dính đầy mồm mép rồi về Suzuno lại cằn nhằn chị mày.''

``Em cảm ơn Coco-oneechan!'' Naomi cười tít mắt, cắn một miếng đường ngọt lịm rồi đưa lên trước mặt Coco: ``Chị có muốn ăn thử không ạ? Ngọt lắm đó!''

Coco ngượng ngùng liếc quanh, thấy không ai để ý mới cúi xuống cắn một miếng nhỏ: ``Ừm\dots\ cũng tàm tạm.''

Khi cả hai ra tới bờ đê cỏ xanh mướt ven sông, gió chiều thổi lồng lộng làm tung bay mái tóc cam của Coco và tà váy trắng của Naomi. Thấy cô bé đi một lúc có vẻ mỏi chân, nàng Rồng chẳng nói chẳng rằng bế bổng cô bé lên rồi đặt ngồi vững chãi trên đôi vai rộng lớn của mình.

``Oa\dots\ cao quá! Nhìn thấy cả cây cầu đằng xa kìa!'' Naomi thích thú reo lên.

Nàng Rồng giữ chặt lấy hai cổ chân nhỏ của cô bé, bật cười sảng khoái: ``Thế nào? Tầm nhìn của Rồng đỉnh hơn người phàm nhiều đúng không?''

Ngồi trên vai Coco, Naomi nhẹ nhàng lấy từ trong giỏ mây ra một chiếc vòng hoa nhỏ được đan khéo léo từ những bông hoa cỏ dại và cúc dại hái ven đường. Bằng động tác cẩn thận nhất, cô bé nhẹ nhàng đặt chiếc vòng hoa lên mái tóc cam bồng bềnh của Coco.

Cô nàng Thế hệ thứ Tư sững người, bước chân dừng lại giữa triền đê đầy gió: ``Cái này\dots\ là gì đây?''

Naomi cúi đầu xuống, hai tay ôm nhẹ lấy trán Coco, cất giọng thủ thỉ chân thành: ``Chiếc vương miện tặng Coco-oneechan. Naomi biết chị đã phải chiến đấu rất nhiều, rất dũng cảm ở thế giới ngoài kia. Nhưng khi về tới ký túc xá nhà mình rồi, chị không cần phải gồng mình làm một con rồng hung dữ hay mạnh mẽ mọi lúc đâu. Ở nhà, chị chỉ cần là chính chị, là Coco-oneechan ấm áp thôi ạ.''

Lời nói non nớt nhưng chứa đựng sự thấu hiểu sâu sắc từ một linh hồn cổ xưa như chạm thẳng vào tận cùng góc khuất trong tim Coco. Suốt bao năm qua, cô luôn là người đi tiên phong, là người gánh vác, là tấm khiên che chở cho đàn em trước mọi mũi dùi dư luận, chưa từng có ai bảo cô rằng cô được phép buông bỏ gánh nặng để sống cho chính mình.

Sống mũi Coco bỗng cay xè. Cô vội vàng ngửa mặt lên nhìn bầu trời hoàng hôn đỏ rực, chớp chớp mắt để ngăn những giọt nước mắt chực trào ra: ``Đúng là\dots\ con nít ranh mà nói chuyện như bà cụ non vậy không biết\dots''

Nhưng bàn tay cô lại siết chặt lấy bàn chân nhỏ của Naomi hơn, ngập tràn sự trân trọng và biết ơn vô hạn.

\ct

Trời sẩm tối, khi Kuon và Kaigou vừa mở cửa bước vào nhà, mùi thơm nồng nàn của bơ nướng và hương trà hoa cúc ngào ngạt đã lan tỏa khắp hành lang.

``Bọn anh/chị về rồi đây!'' cả hai đồng thanh cất tiếng.

``Mừng hai người về nhà, Onii-sama, Onee-sama!'' Suzuno từ trong bếp tươi cười bước ra, trên tay bê khay trà nóng hổi. ``Mọi người mau vào phòng khách đi, có một bất ngờ lớn đang chờ đấy ạ!''

Người Anh cả tháo giày, tò mò bước vào phòng khách và không khỏi ngạc nhiên. Căn phòng đã được dọn dẹp ấm cúng, chiếc bàn thấp ở trung tâm bày đầy những đĩa bánh quy tạo hình dễ thương, hoa quả tươi và những tách trà nghi ngút khói.

Trên chiếc ghế sofa, Ryuuhou vừa mới tỉnh ngủ đang dụi mắt ngơ ngác, trong khi Coco - vẫn đội nguyên chiếc vòng hoa dại trên đầu - đang ngồi uống trà với vẻ mặt cực kỳ hiền lành khác lạ. Bên cạnh cô là Aloe, gương mặt rạng rỡ cầm theo cây đàn ukulele.

Naomi đứng ở giữa phòng, hai tay cầm những tấm thiệp màu sắc vừa hoàn thành ban chiều.

``Hôm nay, em và Suzuno-oneechan mở bữa tiệc trà này để cảm ơn tất cả mọi người,'' Naomi cất giọng dõng dạc, từng bước tiến lại trao thiệp cho từng thành viên.

``Tấm này của Onii-chan, cảm ơn Onii-chan đã vất vả đi làm và luôn bảo vệ cả nhà.''

``Tấm này của Kaigou-oneechan, chúc chị luôn mạnh mẽ và không sợ làm vỡ đồ đạc nữa.''

``Tấm này của Ryuuhou-oneechan, cảm ơn chị vì luôn làm cả nhà cười.''

Và cuối cùng, cô bé trao hai tấm thiệp to nhất, lấp lánh nhất cho Coco và Aloe: ``Và hai tấm này dành cho Coco-oneechan và Aloe-oneechan. Cảm ơn hai chị vì đã đến và trở thành người một nhà với chúng em!''

Mở tấm thiệp ra, bên trong là nét vẽ ngộ nghĩnh hình tám thành viên đang nắm tay nhau dưới mái nhà Hikawa, kèm dòng chữ nắn nót: ``Chào mừng các anh chị về nhà!''

Kaigou nhìn tấm thiệp mà khóe mắt cay cay, cô nàng khẽ chạm vào chiếc nơ cầu vồng trên tóc rồi dang rộng hai tay ôm chầm lấy Naomi: ``Con bé này\dots\ sao mà đáng yêu thế không biết!''

Ryuuhou nhìn sang Coco, lập tức nhếch môi trêu chọc với phong cách độc miệng quen thuộc: ``Úi chà chà, nhìn bà trùm Yakuza đội vòng hoa kìa! Mặt đỏ bừng như tôm luộc rồi kìa Coco-san! Đừng bảo là chị sắp khóc nhè đấy nhé!''

Coco trừng mắt nhưng miệng lại không giấu nổi nụ cười hạnh phúc, liền với tay bốc một chiếc bánh quy nhét thẳng vào mồm Ryuuhou: ``Ăn bánh đi con nhóc nhiều chuyện! Ngon thế này mà không ngậm mồm lại được à!''

Aloe ôm tấm thiệp vào lòng, mắt lấp lánh niềm vui, liền cầm đàn lên gảy một khúc nhạc rộn rã. Cả căn phòng tràn ngập tiếng cười đùa, tiếng trêu chọc chí chóe của Kaigou và Ryuuhou, tiếng trầm trồ khen bánh của Kuon và giọng hát ngọt ngào của Aloe hòa cùng tiếng chuông leng keng từ tay Suzuno.

Từ chiếc đồng hồ trên cổ tay Kuon, giọng nói trầm ấm của Game khẽ vang lên qua loa nhỏ: ``Quét phân tích dữ liệu sinh trắc học hoàn tất: Mức độ cộng hưởng cảm xúc của toàn bộ cư dân ký túc xá đã đạt mức tối đa 100\%. Kết luận: Nơi này không còn là trạm trú ẩn tạm thời, mà là một gia đình vĩnh cửu.''

Kuon nhấp một ngụm trà ấm, nhìn Naomi đang cười rạng rỡ trong vòng tay của các chị gái. Cậu khẽ mỉm cười, thầm cảm ơn số phận đã mang thiên thần nhỏ này đến với cuộc đời mình.

Chính sự ngây thơ và tình yêu thương thuần khiết của Naomi đã trở thành chiếc mỏ neo vững chắc nhất, gắn kết những linh hồn từng chịu nhiều giông bão lại bên nhau dưới một mái ấm bình yên.

%==========================================TRUYỆN PHỤ=========================================
\omk

Khi màn đêm đã buông sâu, tiếng cười đùa trong phòng khách cũng tắt lịm theo sau bữa tiệc trà. Từng người một lần lượt tản về phòng nghỉ, chỉ còn những ngọn đèn ngủ yếu ớt thắp sáng lẻ loi trong căn nhà. Một mùi hương hoa cúc còn vương lại trong không khí, như lưu giữ khoảnh khắc ấm áp vừa trôi qua.

Bước ra ban công để đón lấy làn gió đêm mát lành, Kuon bất ngờ thấy Coco vẫn đứng đó, bóng dáng in dài dưới ánh trăng non. Trong tay nàng Rồng là chiếc vòng hoa dại Naomi tặng lúc chiều, ngón tay khẽ vuốt nhẹ trên những cánh hoa vẫn còn tươi. Cái vẻ ngạo nghễ thường ngày chẳng còn đâu, thay vào đó là một sự lặng thầm hiếm hoi.

``Chưa ngủ à, Coco-chan?'' Kuon bước lại gần, tựa tay vào lan can lạnh giá.

Coco khẽ quay người, ánh trăng rọi lên mái tóc cam bồng bềnh để lại những vệt sáng mờ ảo. Nàng nâng chiếc vòng hoa lên ngang ngực, giọng nói đượm một nỗi niềm khó tả mà trước đây chưa từng có. ``Kuon này\dots\ tôi từng nghĩ sau khi bước ra khỏi ánh đèn sân khấu, cuộc đời mình sẽ chỉ toàn là những chuỗi ngày trống rỗng và đề phòng.''

Cô hướng ánh mắt vào trong nhà, nơi ánh đèn ngủ ấm áp hắt ra từ căn phòng của Naomi và Aloe. Tiếng thở đều đặn của hai cô bé có thể nghe thấy lờ mờ qua khe cửa hé mở, tạo nên một khung cảnh bình yên đến lạ. ``Nhưng con bé Naomi\dots\ và cả cái gia đình kỳ quặc này của cậu\dots\ đã cho tôi cảm giác được sống lại một lần nữa. Cảm ơn cậu vì đã cho tôi một chốn để trở về.''

Kuon bật cười khẽ, vỗ nhẹ lên vai rộng lớn của nàng Rồng như một lời an ủi thầm lặng. Cậu ngước nhìn bầu trời đầy sao, cảm nhận sự bình yên đang lan tỏa trong từng ngóc ngách của ngôi nhà. ``Đó là nhờ công của Naomi đấy. Ở đây, cô không cần phải làm anh hùng của ai cả. Cứ thoải mái là chính mình thôi.''

Dưới ánh trăng rằm vằng vặc phủ lên mái nhà gỗ cổ, những vết thương lòng tưởng chừng không thể lành đã thực sự khép miệng. Những nỗi đau từng đeo bám bấy lâu nay tan biến đi trong làn gió đêm, nhường chỗ cho những hy vọng mới bắt đầu đơm hoa nở.

\end{document}